\documentclass[11pt]{article}
\usepackage[margin=1in]{geometry}
\usepackage{amsmath,booktabs,graphicx}
\usepackage[hidelinks]{hyperref}
\begin{document}
\begin{center}
{\Large\bfseries Supplementary Analysis for PALACE\par}
\smallskip
Auditing Hidden Reasoning Charges in Commercial LLM APIs through Token-Usage Prediction
\end{center}
\appendix
\section{Prediction metrics}
\label{app:metrics}
Let \(y_i>0\) be the reference number of hidden reasoning tokens for request \(i\) and let \(\hat y_i\) be the greedy point prediction. For \(n\) requests with valid predictions, the relative metrics used in this supplement are
\[
\begin{aligned}
\mathrm{AVG} &= \frac{100}{n}\sum_{i=1}^n \frac{|\hat y_i-y_i|}{y_i},\\
\mathrm{AGG} &= 100\frac{\big|\sum_{i=1}^n(\hat y_i-y_i)\big|}{\sum_{i=1}^n y_i}.
\end{aligned}
\]
Both are percentages. AVG weights requests equally. AGG measures the magnitude of net error as a percentage of total reference usage, so underestimates and overestimates can cancel. The signed aggregate error uses the same denominator and retains the sign of the sum of residuals. A negative signed error indicates underprediction.

At tolerance \(\epsilon\), greedy Pass@1 is the percentage of requests whose predictions satisfy \(|\hat y_i-y_i|/y_i<\epsilon\). Pass@5 credits a request if at least one of five sampled predictions meets the criterion. Both metrics require access to the reference count at evaluation time. Figure~3 plots Pass@1 across prediction tolerances for PALACE and LoRA and detection rates across injection rates for CoIn.

\section{Simulated overreporting and audit decisions}
\label{app:protocol}
For a fixed prompt, final answer, reference count \(y_i\), and greedy prediction \(\hat y_i\), the clean report is \(r_i(0)=y_i\). At inflation rate \(\alpha\) the simulated report is \(r_i(\alpha)=\operatorname{round}((1+\alpha)y_i)\). The request score is \(r_i(\alpha)/\max(1,\hat y_i)\). The batch score is the ratio of summed reported counts to summed predictions. We set a separate threshold for each evaluated batch size at the finite-sample 95th percentile of clean calibration scores. A request or batch is flagged when its score exceeds the threshold.

At each tested inflation rate, the true-positive rate (TPR) is the fraction of inflated reports or batches flagged. The false-negative rate is \(1-\mathrm{TPR}\). The false-positive rate (FPR) is the fraction of clean reports or batches flagged at the same threshold. Section~E reports the observed rates on held-out data.

\section{Bounds implied by the reported prediction results}
\label{app:bounds}
For this analytical comparison, use the fixed rule \(r_i>(1+\tau)\hat y_i\). Suppose a prediction passes the \(\epsilon\)-tolerance test, so \((1-\epsilon)y_i\le\hat y_i\le(1+\epsilon)y_i\). Every such clean request remains unflagged if \((1+\tau)(1-\epsilon)\ge1\). Every such inflated report is flagged if \(1+\alpha>(1+\tau)(1+\epsilon)\). With the manuscript's \(\epsilon=0.33\), a fixed \(\tau=0.50\) satisfies both inequalities at \(\alpha=1.00\) because \(1.50(0.67)=1.005\) and \(1.50(1.33)=1.995<2.00\). Under this rule, the reported Pass@1 gives a lower bound on TPR and its complement gives an upper bound on FPR at this operating point, provided that Pass@1 uses the same set of requests and counts invalid predictions as failures. The inequalities remain valid if the published pass test includes equality.

\begin{table}[ht]
\centering
\caption{Bounds derived from the manuscript's reported Pass@1 for the Qwen2.5-1.5B auditor. TPR refers to a simulated doubling of the reference count and FPR to the paired clean reports. The values inherit the rounding of the published percentages.}
\begin{tabular}{lrrrr}
\toprule
 & \multicolumn{2}{c}{PALACE} & \multicolumn{2}{c}{LoRA}\\
\cmidrule(lr){2-3}\cmidrule(lr){4-5}
Domain & TPR floor (\%) & FPR ceiling (\%) & TPR floor (\%) & FPR ceiling (\%)\\
\midrule
General & 88.40 & 11.60 & 85.66 & 14.34\\
Math    & 63.81 & 36.19 & 49.26 & 50.74\\
Coding  & 58.39 & 41.61 & 35.81 & 64.19\\
Medical & 59.71 & 40.29 & 51.93 & 48.07\\
\bottomrule
\end{tabular}
\end{table}

These bounds describe the fixed rule at 100\% inflation. Section~E evaluates the calibrated rule directly at smaller inflation rates.

\section{Reference data, counting, and uncertainty}
\label{app:calibration}
A reference record contains the prompt, the final answer, and a trusted count of the hidden reasoning tokens generated on that request. It should also identify the target model/version, relevant settings such as the reasoning budget and tool policy, and the counting convention. A provider could release these counts without disclosing the traces. Training the estimator requires labeled records. Calibrating its decision threshold requires a held-out set of clean records representative of the workload, including its domain mix and reasoning-length distribution. Calibration should be checked when model version, reasoning policy, or workload changes.

The original benchmark sources include traces from DeepSeek-R1 and DeepSeek-R1-Distill-Llama-70B. The available general-domain preprocessing script extracts the first \texttt{<think>} block and counts its tokens with the \texttt{deepseek-ai/DeepSeek-V3-0324} tokenizer. It tokenizes the block text without the delimiters or an added chat template. The revision evaluates against the released reference counts for all four domains. For API deployment, reference counts follow the service's definition of billable reasoning tokens.

Figure~4 reports the training reward for a Qwen2.5-3B Math auditor trained on 90,000 examples over one epoch. Section~E reports a separate training-size ablation for matched SFT.

We quantify test-sample uncertainty by jointly resampling paired predictions by question and keeping multiple traces from the same question together. These intervals describe variation across samples at fixed model weights. Independently trained SFT runs quantify variation across seeds.

\section{Revision experiments}

\paragraph{Evaluation protocol.}
Each released domain pool of 1,000 examples was split by question hash into 500 calibration and 500 test examples. Newly trained SFT models used separate training-derived validation and calibration splits, with matching test questions removed and hyperparameters chosen on validation. We use greedy predictions and report 95\% intervals from 2,000 question-cluster bootstrap resamples. Parse failures count as misses for Pass@33 and use the training-domain median for other metrics. No parse failures occurred in these runs.

\paragraph{Point prediction and matched baselines.}
Pass@33 gives the percentage of greedy predictions within the 33\% relative-error tolerance. All entries below use 500 held-out examples per domain. The LightGBM models use length and surface features. The MLP uses input and output lengths. Matched SFT fine-tunes all Qwen2.5-1.5B-Instruct parameters for three epochs with answer-token loss and an empty reasoning segment. The selected learning rate was \(10^{-5}\), tied on validation Pass@33 and selected by lower validation MAPE. The tie-break was specified after observing the tie.

\begin{center}
\small
\begin{tabular}{lrrrrr}
\toprule
Method & General & Math & Code & Medical & Pooled\\
\midrule
PALACE domain 1.5B &88.6&64.8&58.2&62.8&68.6\\
PALACE merged 1.5B &87.2&57.2&51.0&60.8&64.0\\
PALACE merged 3B &87.0&64.4&60.2&58.8&67.6\\
Matched SFT, seed 0 &90.4&69.6&66.0&64.2&72.5\\
Matched SFT, seed 1 &92.2&68.0&65.6&65.0&72.7\\
Matched SFT, seed 2 &90.4&67.2&64.2&63.8&71.4\\
LightGBM, per domain &81.2&45.0&46.2&46.8&54.8\\
LightGBM, pooled training &79.8&44.4&48.0&40.2&53.1\\
Two-length MLP &80.8&33.8&48.8&33.0&49.1\\
Domain median &76.2&30.8&27.6&38.4&43.2\\
\bottomrule
\end{tabular}
\end{center}
Pooled PALACE Pass@33 was 68.6\% (95\% CI [66.5, 70.7]). Matched SFT obtained \(72.2\pm0.7\%\) across three seeds. Relative to SFT seeds 0 and 1, PALACE trailed by 3.9 points (CI [2.0, 5.9]) and 4.1 points (CI [2.2, 6.0]), respectively. Pooled MAPE was 28.7\% for domain-specific PALACE, 30.1\% for merged 3B, 32.0\% for merged 1.5B, 41.8\% for per-domain LightGBM, and approximately 25.4\% for matched SFT.

\paragraph{Direct detection.}
For each reference count \(y\), we simulate the report \(\operatorname{round}((1+\alpha)y)\). The request score is the report divided by \(\max(\hat y,1)\). The batch score is the ratio of total reported to total predicted counts. We set size-specific thresholds at the finite-sample 95th percentile of the corresponding clean calibration scores and average over 50 repeated partitions. A flag is a score above the threshold. False-positive rates in the tables are observed rates on held-out clean batches. Section~C provides a separate bound for a different rule.

\begin{center}
\small
\begin{tabular}{rrrrr}
\toprule
Inflation & General & Math & Code & Medical\\
\midrule
10\% &12.5&8.3&7.2&6.9\\
20\% &21.8&11.5&10.9&9.0\\
30\% &36.8&16.7&15.4&11.3\\
50\% &66.9&27.8&24.7&15.7\\
\midrule
\multicolumn{5}{l}{Other methods at 30\% inflation}\\
Matched SFT, seed 0 &34.9&16.7&11.3&10.2\\
PALACE merged 3B &30.5&14.4&12.9&12.3\\
LightGBM, per domain &20.6&10.9&9.0&9.0\\
\bottomrule
\end{tabular}

\smallskip
Single-request TPR (\%) at a nominal 5\% FPR. The first four rows are for domain-specific PALACE.
\end{center}
At 30\% inflation, the pooled ROC-AUC was 0.67 (General: 0.82). Flagging reports above \(1.33\hat y\) detected 47.6\% of 30\%-inflated requests and flagged 21.8\% of clean requests.

\begin{center}
\small
\begin{tabular}{llrrrr}
\toprule
Method & Inflation & \(n=10\) & \(n=50\) & \(n=100\) & \(n=500\)\\
\midrule
PALACE domain 1.5B &10\%&11.9&30.0&50.5&96.5\\
PALACE domain 1.5B &20\%&22.1&69.9&92.6&100\\
PALACE domain 1.5B &30\%&35.5&92.9&99.7&100\\
PALACE merged 3B &30\%&36.6&95.4&99.9&100\\
Matched SFT, seed 0 &30\%&37.5&96.8&100&100\\
LightGBM, per domain &30\%&28.5&83.0&98.0&100\\
PALACE domain 1.5B FPR &0\%&5.0&4.9&5.2&8.4\\
\bottomrule
\end{tabular}

\smallskip
Batch TPR and PALACE clean FPR (\%) for uniformly mixed four-domain batches.
\end{center}
At 30\% inflation, the first tested size exceeding 90\% TPR was 50 for the uniform mixture and for General, Math, and Code individually, and 100 for Medical. At 10\% inflation, 500 was the first tested size exceeding 90\% TPR. Its observed FPR was 8.4\%. The fixed-size calibration pool provides fewer independent windows at this batch size, which makes the tail threshold less precise. Domain-specific PALACE had approximately \(-12\%\) signed batch bias, which persisted as batch size increased and was absorbed by matched calibration.

\paragraph{Calibration sensitivity and intervals.}
For 100-request batches, using uniform calibration for a workload with 70\% Code raised domain-specific PALACE FPR to 25.9\%. Multiplying each prediction in domain \(d\) by \(c_d=\sum_{\mathrm{cal},d} y/\sum_{\mathrm{cal},d}\hat y\) reduced FPR to 3.5\%. After correction, TPR for 30\% inflation in 50-request batches was 94.5\%. The following table compares all evaluated methods under this workload shift.

\begin{center}
\small
\begin{tabular}{lrrr}
\toprule
& \multicolumn{2}{c}{FPR (\%), \(n=100\)} & TPR (\%), \(n=50\)\\
Method & Before correction & After correction & After correction\\
\midrule
PALACE domain 1.5B &25.9&3.5&94.5\\
PALACE merged 3B &1.8&3.5&96.7\\
Matched SFT, seed 0 &3.1&3.1&97.8\\
LightGBM, per domain &1.3&2.8&86.2\\
\bottomrule
\end{tabular}
\end{center}

Varying clean calibration size from 50 to 500 gave mean FPRs of 3--6\%, with greater variability at small sizes. Calibrating only on short traces raised FPR to 19--31\%. Inflating calibration labels by 10\% substantially impaired detection of inflation up to 10\%. In Medical, inflating the SFT training labels by 10\% reduced Pass@33 from 64.2\% to 59.0\% and raised MAPE from 29.2\% to 33.5\%. Applying the same bias to calibration labels gave a 6.4\% flag rate at 10\% inflation and 5.0\% FPR. Nominal 90\% same-domain conformal intervals covered 89.4--91.4\% of PALACE references and 88.4--91.8\% of SFT references. PALACE upper-to-lower interval ratios ranged from 2.0 to 4.6, and coverage in Code and Medical fell to 83--85\% under calibration shift.

\paragraph{Ablations, training size, and routing.}
Removing explicit numeric length cues in both matched SFT training and test changed Pass@33 by \(+0.4,+0.4,-0.6,+1.6\) points for General, Math, Code, and Medical, respectively, with every confidence interval including zero. Removing prompt and answer text while retaining length cues at PALACE inference reduced Math from 63.8\% to 19.2\%, Code from 59.5\% to 27.3\%, Medical from 60.3\% to 17.5\%, and General from 88.9\% to 55.1\%. Swapping length cues across examples had little effect. The SFT training-size experiment gave the following Pass@33 percentages.

\begin{center}
\small
\begin{tabular}{lrrrr}
\toprule
Domain & 10\% & 25\% & 50\% & 100\%\\
\midrule
General &85.8&88.4&89.4&90.4\\
Math &60.8&65.8&66.2&69.6\\
Code &51.0&58.2&62.2&66.0\\
Medical &45.8&57.4&60.6&64.2\\
\bottomrule
\end{tabular}
\end{center}
The reimplemented TF-IDF/logistic regression router achieved 99.08\% domain classification accuracy. Restricting it to eight words or length alone gave 93.2\% and 67.2\%, respectively. Routing cost 0.2 Pass@33 points relative to oracle routing. Oracle-routed domain models outperformed the merged 1.5B model by 4.55 points (95\% CI [2.85, 6.30]). Sampling five candidates at temperature 0.8 improved oracle Pass@5 by 1--3 points over greedy decoding. The shortest reference was 238 tokens. The 1.5B auditor generated 55--95 tokens per request and evaluated approximately 1,000 requests in 20 seconds on one L40S. The training-size sweep and three-seed comparison use matched SFT. The PALACE evaluations use the released checkpoints. The detection experiments use trace-derived reference counts and simulated reporting changes.
\end{document}
